% C-41 -- results.tex cross-section redundancy condense (redundancy_closing_audit_consolidation_2026-09-02.md, P-11)
% Audit copy only -- live file is sections/results.tex.
%
% "it builds no map, retains no spatial memory, and computes no global re-route when blocked"
% appeared near-verbatim in methodology.tex's Gazebo Simulation paragraph (establishing mention)
% and again in results.tex's Reactive Inspection Task discussion -- same pattern R2-01 already
% established for other cross-section repeats (establishing mention + cross-reference, not full
% restatement). Confirmed by Samuel and Diego independently. Condensed the results.tex occurrence
% to a cross-reference, preserving the rest of the paragraph (the BG/WTA mechanism sentence, which
% adds information not in methodology.tex).

% --- results.tex, ~line 558 ---
\revblue{As described in Section~\ref{sec:methodology}, the system's behavior in this task is purely reactive --- no map, no spatial memory, no deliberate re-routing when blocked. What it implements instead is reactive re-routing driven by the basal-ganglia/Winner-Take-All (BG/WTA) arbitration network, which integrates LiDAR and camera information in real time to laterally evade each obstacle while continuing toward the target once visible. Coverage rate --- the percentage of a predefined inspection area visited --- is not reported here, as it presupposes a coverage planner and a known target area that this reactive architecture neither has nor requires (see Section~\ref{ssec:Limitations}). Instead, task-completion indicators that are both measurable with the existing instrumentation and consistent with the system's actual capabilities are used: success rate (SR), final approach precision ($d_{final}$), and time to first visual acquisition of the target ($T_{acquisition}$).}
